\chapter{Propuesta de solución, Metodología y Herramientas}
Este capítulo presenta la propuesta de solución, la metodología y las herramientas utilizadas durante el desarrollo.
\section{Propuesta de solución}
% Explicar la idea de MiApp sin detalles técnicos
La propuesta consiste en desarrollar TransDrakkar, una aplicación web para coordinar el transporte de mobiliario. El cliente solicita la recogida y consulta sus estados; administración revisa solicitudes, gestiona recursos y organiza rutas desde Yecla; el conductor consulta su itinerario y registra entregas o incidencias. El acceso, las sesiones y los permisos constituyen la base del primer incremento y el proceso logístico completa el segundo. Correo y Google disponen de implementación y pruebas simuladas, separadas de sus comprobaciones reales.
\section{Metodología}
% Justificar la elección
% Marco de trabajo Scrum
La organización prevista en la memoria utiliza sprints y toma Scrum como marco de referencia. La división del trabajo en incrementos permite revisar el prototipo y ajustar las siguientes tareas. DoneTonic se utiliza como soporte para recoger la planificación y el seguimiento; las evidencias concretas de cada sprint se incorporan al Anexo A. Las fechas, estimaciones y horas deben coincidir con sus registros reales.
En el repositorio se observa un flujo de trabajo con ramas y solicitudes de integración. GitHub Actions define comprobaciones automatizadas del código y del contenedor. Este flujo permite revisar cambios antes de integrarlos en la rama principal. Los hitos técnicos del repositorio no se identifican automáticamente con sprints: su correspondencia se documentará a partir de DoneTonic.
\subsection{Seguimiento con DoneTonic}
El backlog permite relacionar el trabajo previsto con las tareas de cada sprint. Para documentar el proceso se incorporarán capturas del backlog y del tablero, explicando qué alcance y qué estado representan. Cada captura debe identificar el sprint y el momento en el que se obtuvo; una captura del estado final no permite reconstruir por sí sola la evolución diaria.
\begin{figure}[htbp]
\centering
\includegraphics[width=0.95\textwidth]{figs/donetonic-backlog.jpg}
\caption{Vista parcial del backlog ordenado por prioridad en DoneTonic, capturada el 27 de octubre de 2026. Se muestran PBI de pruebas, persistencia, configuración y permisos.}
\label{fig:donetonic-backlog}
\end{figure}
\section{Herramientas}
% Gestión del proyecto
\subsection{Gestión del proyecto y versiones}
DoneTonic sirve como soporte de planificación y seguimiento. Git y GitHub mantienen el historial del código, las ramas y las solicitudes de integración. GitHub Actions define la integración continua con pruebas automatizadas.
% Desarrollo (con las librerias)
\subsection{Desarrollo y despliegue}
El cliente está desarrollado con HTML, CSS y JavaScript. El servidor utiliza Node.js 24 y Express. MongoDB, mediante su controlador para Node.js, proporciona persistencia de cuentas, sesiones y transporte, con transacciones para cambios relacionados y bcrypt protege las contraseñas. Nodemailer integra el envío SMTP y el acceso con Google utiliza OAuth con PKCE. SQLite y el ejecutor de pruebas de Node.js permiten realizar pruebas aisladas. Docker empaqueta la aplicación; Cloud Build y Cloud Run proporcionan la configuración de construcción y despliegue, mientras que Atlas aloja la base de datos.
% Documentación (Overleaf...)
\subsection{Documentación}
La memoria se redacta en LaTeX mediante Overleaf. El README describe ejecución, API y pruebas; la guía de despliegue recoge la configuración operativa. Las imágenes de la memoria se guardan en la carpeta \texttt{figs} y las referencias bibliográficas en \texttt{bib/ref.bib}.
% Comunicación (TFG)
La memoria y las evidencias de seguimiento permiten comunicar el progreso del proyecto. Los canales y reuniones que se hayan utilizado se concretarán a partir del trabajo real, sin atribuir ceremonias o herramientas no registradas.
